\paragraph{Resource preflight and adaptive refinement.}
The scalable workflow first creates a structural partition and then performs
bounded exact-support discovery.  The old sharded planner materializes global
origin support and retains per-column patterns and future construction tasks;
it is therefore a compatibility path for reviewed small cases, not a
full-network preflight.  The streaming preflight instead prepares and analyzes
one destination group at a time.  It reduces exact reachable measurement rows
immediately to group and block summaries, writes an atomic checkpoint, and
releases group-local routing, reachability, pattern, and mapping state before
continuing.  Retained state is proportional to destination and block summaries,
not to OD--measurement support entries.

Four explicit modes distinguish structural inspection, deterministic sampled
exact support, complete streaming exact support, and an explicitly authorized
materialized compatibility plan.  The materialized mode is rejected unless a
prior conservative logical-support estimate fits its retained-state budget.
Sampled selection includes small, median, 95th-percentile, and maximum
structural groups.  Its extrapolations report observed ranges and assumptions;
sampled completion alone cannot authorize a complete-network pilot.

Elapsed time, process RSS, temporary group state, retained summaries, support
rows, nonzeros, and operator bytes are effective limits.  Budget exhaustion or
interruption produces a typed partial result and a valid atomic checkpoint.
Resume validates scenario, assignment, OD layout, fixed demand, measurement
mapping, routing, partition, configuration, and schema fingerprints.  It
neither repeats completed groups nor skips pending groups.

Checkpoint identity separates numerical semantics from invocation policy.
The semantic fingerprint covers mode, deterministic sampling and selected
groups, support chunking and tolerance, persisted representation, and schema;
the authoritative data and routing fingerprints remain mandatory.  Elapsed and
memory allowances, checkpoint cadence, and deterministic worker settings have
a separate recorded policy fingerprint and may change compatibly on resume.
Tightened limits are rejected before work if retained summaries or accepted
blocks already violate them.

Cumulative support time is never reset and includes discarded work from an
incomplete group.  In contrast, the elapsed-time allowance applies only to the
current invocation and its timer begins at zero after each resume.  Reports
give previous, current, and cumulative time, invocation count, allowance
overshoot, and whether stopping occurred before a group or inside a bounded
chunk.  Thus exhausting a 120-second allowance does not cause the next
120-second invocation to stop immediately.  An unfinished group is not
committed and may be recomputed, while every completed group is skipped.

Only after support discovery does the resource preflight select the smallest,
median, 95th-percentile, largest, or explicitly named blocks.  Before invoking
any builder it estimates CSR values and indices, transpose storage,
construction temporaries, retained cache, and solver work arrays.  An unsafe
block is rejected before allocation; no unselected operator is constructed.
For accepted blocks it records machine memory,
logical and physical CPUs, cache storage, operator and local-solver memory,
construction and product times, and cache/checkpoint sizes.  The \texttt{auto},
\texttt{laptop}, \texttt{workstation}, and \texttt{server} profiles apply
different conservative fractions of available memory.  Worker count is bounded
by both CPUs and safety-adjusted measured peak memory:
\[
n_{\mathrm{workers}}
\leq
\min\!\left(
n_{\mathrm{CPU}},
\left\lfloor\frac{M-M_0}{M_B}\right\rfloor
\right),
\]
where $M$ is the selected job budget, $M_0$ is coordinator and assignment
memory, and $M_B$ is the conservative measured worker peak.  Estimated cache
storage is checked separately.  The structured recommendation contains hard
variable and nonzero ceilings, worker and thread allocation, peak memory,
cache size, first-sweep and cache-hit sweep times, and uncertainty.  Automatic
sizing is never activated until the exact recommendation fingerprint is
explicitly accepted.

Selected fixed-routing blocks are constructed directly from their authoritative
assignment inputs, compact free-to-active OD layout, destination-group routing,
measurement mapping, and persisted exact support.  The production builder does
not require a complete measurement operator.  It combines one or more base OD
chunks into a separately configured, memory-guarded OD batch. For each batch,
one JIT-compiled forward dynamic program computes node reach probabilities once;
fixed-size mapped-edge gathers
then aggregate only bounded supported-measurement chunks on the host. Numerical
zeros are filtered before canonical COO-to-CSR assembly.  Both CSR and
CSC storage retain only the union of supported rows, while forward products
scatter into the complete measurement coordinate system and transpose products
gather from it.  Consequently, the forbidden dense array of shape
$n_{\mathcal M}\times |B|$ is never allocated: the largest dense numerical
buffers are bounded by the effective OD batch times the node count, mapped-edge chunk,
or supported-measurement chunk.

Measurement-row lookup, enabled-link filtering, global-to-local row translation,
and padded mapped-edge plans are invariant for the block and are prepared once,
outside the OD-batch loop. The requested batch factor is reduced automatically
until conservative reach, mapping, sparse-assembly, routing, solver, and retained
storage estimates satisfy both temporary and worker ceilings. Failure of even
one base batch is reported before numerical allocation. Because batching does
not change the canonical per-column accumulation order, safe batch sizes share
one logical cache identity.

Exact selected support and each numerical block are persisted independently by
atomic replacement.  Their identities include package and schema versions,
all scenario, assignment, OD-layout, fixed-demand, measurement, routing, and
partition fingerprints, the block and support fingerprints, fixed routing
parameter, tolerance, dtype, and canonical accumulation contract. Pure
performance chunk sizes are excluded from the logical cache identity. Content
hashes, sparse dimensions, row coordinates, and dtype are validated before
reuse; corrupt numerical caches are discarded and rebuilt, while corrupt or
incompatible support artifacts are rejected.  A fresh process can therefore
load a block without repeating routing or numerical construction.  In-memory
retention is explicitly bounded and releasable.

The accepted measurements also define a conservative memory/runtime cost model.
Before estimation, any unsafe candidate block is bisected deterministically
until every child satisfies the accepted limits.  Coverage remains exact,
children retain a conservative measurement support, and blocks are never
silently enlarged.  If a minimum-size block is still unsafe, a resource guard
stops the workflow before its operator is constructed.  Refinement creates a
new partition fingerprint, and the revised schedule is written in the initial
checkpoint.  A checkpoint for the obsolete parent partition therefore cannot
resume the refined run.

This block-coordinate implementation is presently restricted to the
fixed-routing linear MAP objective.  It must not be used when routing changes
with the OD iterate or when the required conditional prior cannot be evaluated
exactly.

